\subsection{Proof of \Cref{lem:cpfuse-kkt}}

For brevity, write
\[
\gamma_i
=
\log\left(
\frac{\model(y_{<t}\mid x)}
     {\modeli(y_{<t}\mid x)}
\right).
\]
The optimization problem in \Cref{eq:cpfuse-opt-problem} can then be written as
\begin{equation}
\begin{aligned}
\min_{p^*,\,\tau} \quad & \tau \\
\text{s.t.} \quad &
\KL\bigl(p^* \,\|\, \modeli(\cdot \mid y_{<t},x)\bigr)
+ \gamma_i
\le \tau,
\qquad \forall i,
\end{aligned}
\end{equation}
with \(p^*\) constrained to lie on the simplex over \(V\).

Let \(\lambda_i \ge 0\) be the multipliers of the two inequality constraints and let \(\mu\) enforce simplex normalization. The KKT conditions give
\[
\sum_i \lambda_i = 1
\]
from stationarity in \(\tau\), and for every coordinate with \(p^*(y_t) > 0\),
\begin{equation}
\sum_i \lambda_i
\Bigl(
\log p^*(y_t)
- \log \modeli(y_t \mid y_{<t}, x)
+ 1
\Bigr)
+ \mu
= 0.
\end{equation}
Using \(\sum_i \lambda_i = 1\), this becomes
\begin{equation}
\begin{aligned}
\log p^*(y_t)
= {} &
\lambda_1 \log \modelone(y_t \mid y_{<t},x) \\
& + \lambda_2 \log \modeltwo(y_t \mid y_{<t},x)
+ \gamma_t,
\end{aligned}
\end{equation}
where \(\gamma_t = -1-\mu\). Setting \(\alpha_t = \lambda_1\) and \(\beta_t = \lambda_2\) yields the claimed representation. For coordinates with \(p^*(y_t)=0\), the same identity is read in the extended-real convention \(\log(0)=-\infty\).

\subsection{Proof of \Cref{lem:cpfuse-balance}}

Assume that both \(\modelone\) and \(\modeltwo\) have full support, and define
\[
\gamma_i
=
\log\left(
\frac{\model(y_{<t}\mid x)}
     {\modeli(y_{<t}\mid x)}
\right).
\]
Also define
\[
f_i(p)
=
\KL\bigl(p \,\|\, \modeli(\cdot \mid y_{<t},x)\bigr)
+
\gamma_i.
\]
The optimization problem in \Cref{eq:cpfuse-opt-problem} is therefore
\begin{equation}
\min_{p} \max\{f_1(p),f_2(p)\}.
\end{equation}
At any optimizer \(p^*\), either both terms are active, or exactly one of them is active.

\begin{enumerate}
    \item Both constraints from \Cref{eq:cpfuse-opt-problem} are tight. In that case,
    \begin{equation}
    \begin{aligned}
    \KL\bigl(p^* \,\|\, \modelone(\cdot \mid y_{<t},x)\bigr) + \gamma_1
    &=
    \KL\bigl(p^* \,\|\, \modeltwo(\cdot \mid y_{<t},x)\bigr) + \gamma_2.
    \end{aligned}
    \end{equation}
    Expanding the KL terms shows that this is equivalent to
    \begin{equation}
    \begin{aligned}
    \mathbb{E}_{y_t \sim p^*}\log \modelone(y_{\le t}\mid x)
    &=
    \mathbb{E}_{y_t \sim p^*}\log \modeltwo(y_{\le t}\mid x),
    \end{aligned}
    \end{equation}
    which is exactly condition (1) in \Cref{lem:cpfuse-balance}.

    \item Exactly one term is active. Then \(p^*\) minimizes that active term. Since each function \(f_i\) is minimized on the simplex at the corresponding base conditional, the only possibilities are that \(p^*\) coincides with one of the two base conditionals.

    We now rule out the first possibility. Suppose by contradiction that \(p^* = \modelone(\cdot \mid y_{<t},x)\). Then
    \begin{align}
    \KL\bigl(p^* \,\|\, \modeltwo(\cdot \mid y_{<t},x)\bigr) + \gamma_2
    &>
    \KL\bigl(\modeltwo(\cdot \mid y_{<t},x) \,\|\, \modeltwo(\cdot \mid y_{<t},x)\bigr) + \gamma_2 \\
    &=
    \gamma_2 \\
    &>
    \gamma_1 \\
    &=
    \KL\bigl(p^* \,\|\, \modelone(\cdot \mid y_{<t},x)\bigr) + \gamma_1.
    \end{align}
    Here the strict inequality \(\gamma_2>\gamma_1\) follows from the assumption \(\modelone(y_{<t}\mid x)>\modeltwo(y_{<t}\mid x)\). Hence the objective at \(p^*=\modelone(\cdot \mid y_{<t},x)\) is controlled by the second term, so \(p^*\) cannot be optimal. Therefore the only remaining one-active case is
    \begin{equation}
    p^* = \modeltwo(\cdot \mid y_{<t},x),
    \end{equation}
    which is exactly condition (2) in \Cref{lem:cpfuse-balance}.
\end{enumerate}

Finally, if one of the base conditionals assigns zero probability to some token \(y_t\), then any optimizer with finite objective must also assign zero probability to that token. Without the full-support assumption, the optimizer need not satisfy either of the two conclusions above in exactly the same form.
